\section*{Hoofdstuk \ref{ch:b3:renal-osmoregulation} --- Nierfysiologie en osmoregulatie}

\begin{solution}{exo:b3:renal-osmoregulation:1}
\omterm{def:b3:renal-osmoregulation:nephron}{Glomerulus}: filtreert plasma op de eiwitten na. \omterm{def:b3:renal-osmoregulation:nephron}{Proximale tubulus}:
neemt twee derde van het zout en het water terug, alle glucose en
aminozuren, bicarbonaat. Dalende tak van Henle: staat water af aan het
merg. Stijgende tak: pompt NaCl uit, laat geen water door. Distale
tubulus: stelt het natrium bij (\omterm{def:b3:renal-osmoregulation:controls}{aldosteron}), scheidt kalium uit.
\omterm{def:b3:renal-osmoregulation:nephron}{Verzamelbuis}: terugname van water onder vasopressine, verwerking van
protonen en ureum.
\end{solution}

\begin{solution}{exo:b3:renal-osmoregulation:2}
Klaring: het volume plasma dat per tijdseenheid van de stof wordt
ontdaan, $UV/P$. Hier $100\times 1.5/2 = \qty{75}{mL/min}$, onder de
GFR van \qty{120}{mL/min}: de stof wordt gefiltreerd en deels
teruggenomen (ongeveer \qty{38}{\%} van de gefiltreerde hoeveelheid).
\end{solution}

\begin{solution}{exo:b3:renal-osmoregulation:3}
Een zoetwatervis is zouter dan zijn medium: water komt door osmose door
de kieuwen naar binnen, dus zou drinken de vloed alleen vergroten; hij
scheidt overvloedige verdunde urine uit en zijn chloridecellen in de
kieuwen nemen Na$^{+}$ en Cl$^{-}$ actief uit het water op. Een zeevis
is minder zout dan de zee: hij verliest water door de kieuwen, drinkt
zeewater om het te vervangen, en de chloridecellen pompen het
ingeslikte zout naar buiten.
\end{solution}

\begin{solution}{exo:b3:renal-osmoregulation:4}
Het \omterm{prop:b3:renal-osmoregulation:countercurrent}{enkelvoudige effect} is het verschil van \qty{200}{mOsm/L} dat de
dikke stijgende tak tussen het interstitium en haar eigen vloeistof in
stand houdt door NaCl uit te pompen zonder water te laten volgen. De
stijgende tak moet ondoorlaatbaar voor water zijn, anders zou water het
zout volgen en zou het verschil verdwijnen; de dalende tak moet
waterdoorlatend zijn, zodat de vloeistof die de bocht bereikt al door
het interstitium is geconcentreerd en elke trap op zouter vloeistof
werkt dan de vorige. Twee takken met dezelfde doorlaatbaarheid zouden
niet kunnen vermenigvuldigen.
\end{solution}

\begin{solution}{exo:b3:renal-osmoregulation:5}
Aanvoerend eind: $60 - 18 - 28 = \qty{14}{mmHg}$; afvoerend eind: $60 - 18 -
35 = \qty{7}{mmHg}$. Naarmate er water wordt gefiltreerd raken de
achterblijvende eiwitten geconcentreerd en stijgt de oncotische druk
langs het haarvat; bereikt zij \qty{42}{mmHg}, dan is de netto druk nul
en stopt de filtratie vóór het einde (filtratie-evenwicht). De GFR
hangt dan van de plasmastroom af: een snellere stroom concentreert de
eiwitten trager, zodat de filtratie over een langer stuk doorgaat en de
GFR met de stroom meestijgt.
\end{solution}

\begin{solution}{exo:b3:renal-osmoregulation:6}
De gefiltreerde vracht is $\qty{40}{mL/min}\times\qty{0.14}{mmol/mL} =
\qty{5.6}{mmol/min}$. Uitscheiding: $70\times 1 = \qty{0.07}{mmol/min}$.
Teruggenomen: $5.53/5.6 = \qty{98.75}{\%}$. De natriumklaring is
$70\times 1/140 = \qty{0.5}{mL/min}$.
\end{solution}

\begin{solution}{exo:b3:renal-osmoregulation:7}
$900/1200 = \qty{0.75}{L}$ per dag. Bij \qty{400}{mOsm/L}: $900/400 =
\qty{2.25}{L}$ urine, dus moet er ongeveer $2.25 + 0.9 = \qty{3.15}{L}$
worden gedronken (het water in het voedsel niet meegerekend): de
patiënt bij wie het concentratiemechanisme faalt moet drinken, en
verkeert in gevaar als dat niet lukt.
\end{solution}

\begin{solution}{exo:b3:renal-osmoregulation:8}
(a) Diabetes insipidus: \qtyrange{10}{15}{L} urine bij
\qtyrange{50}{100}{mOsm/L}, een hoog plasmanatrium zodra het drinken
achterblijft, onophoudelijke dorst. (b) Ongepaste vasopressine: karige
geconcentreerde urine, water vastgehouden, plasmanatrium laag door
verdunning, de dorst niet verhoogd hoewel het drinken doorgaat ---
verwardheid en toevallen wanneer het natrium ver zakt. (c) Overmaat aan
\omterm{def:b3:renal-osmoregulation:controls}{aldosteron}: natrium met water vastgehouden, volume uitgezet, bloeddruk
hoog, kalium verloren (laag plasma-K$^{+}$), plasmanatrium bijna
normaal omdat water het zout volgt en de nier deels ontsnapt;
urinevolume normaal. (d) Zoutloos dieet: renine en \omterm{def:b3:renal-osmoregulation:controls}{aldosteron} stijgen
binnen een dag, het urinenatrium daalt tot enkele millimol per dag, het
plasmanatrium blijft op peil, volume en druk iets lager.
\end{solution}

\begin{solution}{exo:b3:renal-osmoregulation:9}
Verdunde urine: osmolaire klaring $100\times 6/300 = \qty{2}{mL/min}$;
\omterm{prop:b3:renal-osmoregulation:water}{vrijwaterklaring} $6 - 2 = +\qty{4}{mL/min}$: de nier loost water zonder
opgeloste stof, zoals na een waterbelasting met vasopressine uitgezet.
Geconcentreerde urine: osmolaire klaring $1000\times 0.6/300 =
\qty{2}{mL/min}$; \omterm{prop:b3:renal-osmoregulation:water}{vrijwaterklaring} $0.6 - 2 = -\qty{1.4}{mL/min}$: er
wordt \qty{1.4}{mL/min} vrij water aan het lichaam teruggegeven ---
uitdroging, vasopressine hoog.
\end{solution}

\begin{solution}{exo:b3:renal-osmoregulation:10}
Nummer de trappen vanaf de schors; zij het interstitium op trap $k$
gelijk aan $I_{k}$. De dalende vloeistof op $k$ is gelijk aan $I_{k}$;
de stijgende vloeistof op $k$ is $I_{k} - 200$; en de stijgende
vloeistof op $k$ is wat de bocht verliet en de trappen
$n, n-1, \dots, k+1$ passeerde. In de stationaire toestand verlaat de
vloeistof die met \qty{300}{} binnenkomt de top van de stijgende tak
met $I_{1} - 200$, en overtreft het interstitium van elke trap dat van
de trap erboven met hoogstens het \omterm{prop:b3:renal-osmoregulation:countercurrent}{enkelvoudige effect}, dus
$I_{k} \le 300 + 200k$ met gelijkheid wanneer elke trap op volle
sterkte werkt: de bocht bereikt $300 + 200n$. Het verval wordt dus
begrensd door de lengte van de lis (het aantal trappen) maal het
\omterm{prop:b3:renal-osmoregulation:countercurrent}{enkelvoudige effect}, niet door het pompvermogen van enige cel. $n$ is
de lengte van de lis gedeeld door de afstand waarover de vloeistof met
het interstitium in evenwicht komt: lange juxtamedullaire lissen, en de
zeer lange lissen van woestijnknaagdieren, geven een grote $n$; het
vermogen van de pomp en het uitspoelen van opgeloste stof door de vasa
recta begrenzen het effect.
\end{solution}

\begin{solution}{exo:b3:renal-osmoregulation:11}
\omterm{def:b3:renal-osmoregulation:comparative}{Metabool water}: $5\times 0.8 = \qty{4}{g}$ zetmeel geeft $4\times 0.6
= \qty{2.4}{g}$. Verliezen: urine $3/5500\ \unit{L} = \qty{0.55}{mL}$;
verdamping \qty{1.5}{g} (al na aftrek van wat de neus terugwint).
Totaal verlies \qty{2.05}{g} tegen een winst van \qty{2.4}{g}: een
overschot van \qty{0.35}{g}, dat het kleine verlies in de ontlasting
dekt, en het hygroscopische water van de zaden is een toegift. Hij
overleeft zonder te drinken. Zonder de terugwinning in de neus zou het
ademverlies $1.5/0.3 = \qty{5}{g}$ zijn, en zou hij binnen enkele dagen
sterven.
\end{solution}

\begin{solution}{exo:b3:renal-osmoregulation:12}
Bij gelijkstroom komen het bloed en het dialysaat halverwege het
membraan in evenwicht en daalt het concentratieverschil dat de diffusie
drijft tot nul: in het beste geval vertrekt het bloed met de helft van
zijn ureum. Bij tegenstroom ontmoet het meest verse dialysaat het best
gereinigde bloed en blijft het verval over het hele membraan behouden,
zodat het bloed bijna met de ingangsconcentratie van het dialysaat kan
vertrekken --- hetzelfde beginsel als de \omterm{def:b3:renal-osmoregulation:nephron}{lis van Henle}. Klaring:
$300\times 0.7 =
\qty{210}{mL/min}$ gedurende $3\times 4 = \qty{12}{h}$ per week, d.w.z.\
$210\times 720 = \qty{151}{L}$ per week; twee nieren klaren bij
\qty{125}{mL/min} $125\times\num{10080} = \qty{1260}{L}$. De machine
doet ongeveer een achtste van het werk van de nieren, over de tijd
gemiddeld \qty{15}{mL/min} --- genoeg om te leven, niet genoeg om
gezond te zijn.
\end{solution}

\begin{solution}{pb:b3:renal-osmoregulation:1}
\textbf{1.} Netto druk $55 - 15 - 30 = \qty{10}{mmHg}$; GFR $= 12.5
\times 10 = \qty{125}{mL/min}$.
\textbf{2.} $125\times 1440 = \qty{180}{L/dag}$; $180/3 = 60$ maal.
\textbf{3.} Inuline: $62.5\times 1/0.5 = \qty{125}{mL/min}$, gelijk aan
de GFR. Creatinine: $1.3\times 1/0.01 = \qty{130}{mL/min}$, iets
groter omdat de \omterm{def:b3:renal-osmoregulation:nephron}{proximale tubulus} naast het gefiltreerde nog een
kleine hoeveelheid creatinine uitscheidt.
\textbf{4.} PAH-klaring $13\times 1/0.02 = \qty{650}{mL/min}$, de
plasmastroom door de nier; bloedstroom $650/(1 - 0.45) = \qty{1180}{mL/min}$,
ongeveer \qty{1.2}{L/min}; filtratiefractie $125/650 = 0.19$.
\textbf{5.} Netto druk \qty{15}{mmHg}, GFR \qty{188}{mL/min} (in
werkelijkheid minder, omdat de oncotische druk langs het haarvat
stijgt). Angiotensine II vernauwt het afvoerende arteriooltje en
verhoogt de druk in de \omterm{def:b3:renal-osmoregulation:nephron}{glomerulus}; het blokkeren ervan verlaagt de druk
en de GFR een weinig. Bij diabetes filtreren de glomeruli onder hoge
druk en slijt het filter; de druk verlagen vertraagt de schade met
jaren.
\textbf{6.} De podocyten en hun spleetdiafragma's (of de lading van het
basaalmembraan): de laag die albumine tegenhoudt. Het verlies van
\qty{3}{g/dag} verlaagt de oncotische druk van het plasma; overal
verlaat vocht de haarvaten en zwellen de weefsels --- het oedeem van
het nefrotisch syndroom.
\textbf{7.} Gefiltreerd $125\times 5 = \qty{0.625}{mmol/min}$;
uitgescheiden $250\times 1 = \qty{0.25}{mmol/min}$; teruggenomen
$0.375$, d.w.z.\ \qty{60}{\%}; klaring $250/5 = \qty{50}{mL/min}$.
\textbf{8.} Gefiltreerd $180\times 0.14 = \qty{25.2}{mol/dag}$ aan Na$^{+}$,
\qty{580}{g} natrium, \qty{1.47}{kg} als NaCl. Uitgescheiden $100\times 1.44
= \qty{144}{mmol/dag}$; teruggenomen $1 - 144/\num{25200} = \qty{99.4}{\%}$.
Er zou bijna anderhalve kilogram zout per dag verloren gaan.
\textbf{9.} Gefiltreerd $125\times 5 = \qty{0.625}{mmol/min}$, d.w.z.\
\qty{112}{mg/min}, \qty{162}{g/dag}; niets uitgescheiden, want dit ligt
onder $T_{m}$. De gefiltreerde vracht bereikt $T_{m}$ bij $2.1/0.125 =
\qty{16.8}{mmol/L}$ (\qty{300}{mg/dL}) --- hoger dan de waargenomen
drempel van ongeveer \qty{10}{mmol/L}, omdat de \omterm{def:b3:renal-osmoregulation:nephron}{nefronen} onderling
verschillen en de zwakste het eerst overlopen (de spreiding van de
titratiekromme).
\textbf{10.} Gefiltreerd $125\times 20 = \qty{2.5}{mmol/min}$;
uitgescheiden $2.5 - 2.1 = \qty{0.4}{mmol/min}$, \qty{72}{mg/min}, $576$
mmol of \qty{104}{g} per dag. Extra urine $576/300 = \qty{1.9}{L}$. De
glucose sleept water de urine in (osmotische diurese), de patiënt
droogt uit en heeft dorst; \qty{104}{g} glucose is \qty{1.8}{MJ} die per
dag verloren gaat, en de cellen, die geen glucose kunnen opnemen,
verbranden vet en eiwit: het gewicht daalt.
\textbf{11.} Teruggenomen $\approx \qty{25}{mol/dag}$; ATP $25/3 =
\qty{8.4}{mol/dag}$; $8.4\times 50 = \qty{420}{kJ}$, ongeveer \qty{6}{\%}
van \qty{7}{MJ}.
\textbf{12.} Om elk afvalproduct uit te scheiden zou de tubulus een
transporteiwit nodig hebben dat het herkent --- duizenden ervan, en
geen enkel voor een molecuul dat nooit eerder is ontmoet. Alles wat
klein is filtreren en de paar honderd soorten terugnemen die het
lichaam wil hebben, vergt alleen transporteiwitten voor wat bekend
staat als waardevol; alles wat niet wordt herkend vertrekt vanzelf. De
prijs is de energie van vraag~11.
\textbf{13.} Osmolaire klaring $600\times 1/290 = \qty{2.07}{mL/min}$;
\omterm{prop:b3:renal-osmoregulation:water}{vrijwaterklaring} $1 - 2.07 = -\qty{1.07}{mL/min}$: de nier spaart
water.
\textbf{14.} Minimum $600/1200 = \qty{0.5}{L}$; maximum $600/50 =
\qty{12}{L}$; dagelijkse behoefte bij het minimum $0.5 + 0.9 = \qty{1.4}{L}$.
\textbf{15.} \qty{1000}{mOsm} bij \qty{1200}{mOsm/L} vergt
\qty{0.83}{L} urine: netto winst \qty{0.17}{L} vóór de opgeloste stof
van de dag. Daarna: de totale opgeloste stof van de dag is
\qty{1600}{mOsm}, \qty{1.33}{L} urine, plus \qty{0.9}{L} onmerkbaar,
\qty{2.23}{L} eruit tegen \qty{1}{L} erin --- een tekort van
\qty{1.23}{L}, tegen \qty{1.4}{L} zonder te drinken. De liter zeewater
verbetert de balans met slechts \qty{0.17}{L}.
\textbf{16.} Urine van \qty{50}{mOsm/L} voor de dag: $600/50 =
\qty{12}{L}$; er moet dagelijks \qty{12.9}{L} worden gedronken, ongeveer \qty{13}{L}.
\textbf{17.} $B$ liter bier brengt $20B$ mOsm; de maximale urine is
$(600 + 20B)/50 = 12 + 0.4B$ liter, en zij moet $B - 0.9$ liter dragen:
$B - 0.9 \le 12 + 0.4B$, $B \le \qty{21.5}{L}$. Ongeveer twintig liter
--- maar als er verder weinig wordt gegeten daalt de opgeloste stof van
de dag tot \qty{200}{mOsm} en zakt de grens tot ongeveer \qty{8}{L}: de
hyponatriëmie van zware drinkers die niet eten.
\textbf{18.} $298/290 - 1 = \qty{2.8}{\%}$: ruim voorbij de \qty{1}{\%}
die vasopressine en dorst uitlokt; beide zijn maximaal. Bij constante
opgeloste stof: $290\times 42 = 298\times V$, $V = \qty{40.9}{L}$, een
tekort van ongeveer \qty{1.1}{L}.
\textbf{19.} Zes liter water verdunt het plasma terwijl het zweet zout
wegvoert; en inspanning, pijn en stress maken vasopressine vrij
ongeacht de \omterm{def:b3:renal-osmoregulation:problem}{osmolariteit}, zodat de nier de urine niet genoeg kan
verdunnen om het overschot uit te scheiden. Het plasmanatrium daalt, de
hersencellen zwellen --- er zijn toevallen en sterfgevallen op gevolgd.
De nier zou vasopressine uitgezet moeten hebben en \qty{6}{L} urine bij
$U_{\min}$ moeten maken; de loper kan het beste op zijn dorst drinken
en het zout aanvullen.
\textbf{20.} $3/5500\ \unit{L} = \qty{0.55}{mL}$ per dag. Een menselijke
nier zou $3/1200 = \qty{2.5}{mL}$ nodig hebben voor dezelfde opgeloste
stof: de kangoeroerat gebruikt een vijfde van het water.
\textbf{21.} Erin: \qty{2.4}{g}. Eruit: $1.6 + 0.3 + 0.55 = \qty{2.45}{g}$.
Tot op \qty{0.05}{g} per dag in evenwicht, wat het hygroscopische water
van de zaden (enkele procenten van \qty{5}{g}) levert.
\textbf{22.} Water $500\times 0.65 = \qty{325}{L}$; een kwart is
\qty{81}{L}; bij \qty{5}{L/dag} zestien dagen.
\textbf{23.} $150/800 = \qty{0.19}{L}$, ongeveer \qty{190}{mL} pekel
--- minder dan de \qty{100}{mL} zeewater en het water in de vis samen,
dus wint de vogel water. Zijn nier zou bij \qty{600}{mOsm/L}
(\qty{300}{mmol/L} NaCl) \qty{0.5}{L} urine nodig hebben om hetzelfde
zout te dragen, meer water dan hij binnenkreeg: een nier die het niet
van zeewater wint kan een vogel het niet laten drinken, en de klier,
bij tweemaal de concentratie van de zee, wel.
\textbf{24.} \qty{30}{g} water erin, dus ongeveer \qty{30}{mL} urine per
dag, bij een klein deel van de \omterm{def:b3:renal-osmoregulation:problem}{osmolariteit} van het bloed (enkele
tientallen mOsm/L tegen \num{300}). Zout verdwijnt in die urine en door
diffusie over de kieuwen; zonder actieve opname door de chloridecellen
uit water dat een millimol per liter bevat, zou de vis binnen dagen
uitgeput raken.
\textbf{25.} GFR \qty{125}{mL/min}, \qty{180}{L} per dag; verplichte
urine \qty{0.5}{L} per dag; een liter zeewater levert netto
\qty{0.17}{L} zodra het zout is uitgescheiden --- bijna niets.
\end{solution}
